\section{RESULT AND DISCUSSION}
\label{sec:results}

We organize our empirical findings across the three core research questions: Section~\ref{subsec:res_rq1} evaluates on-device ViSoBERT against baselines (RQ1); Section~\ref{subsec:res_rq2} ablates the Tier 2 forensic council and single-LLM baselines on hard challenges (RQ2); and Section~\ref{subsec:res_rq3} assesses end-to-end cascade latency, cost, and routing efficiency (RQ3).

\subsection{RQ1: Edge Screener Performance and Significance}
\label{subsec:res_rq1}

Table~\ref{tab:benchmark_summary} presents multi-seed validation performance ($K=5$ random seeds) alongside deterministic evaluation on the locked Frozen Holdout partition ($N=267$: $192$ ham, $75$ scam).

\begin{table*}[t]
\centering
\caption{Multi-Seed Validation Performance across 5 Independent Training Runs ($K=5$) and Frozen Holdout Benchmark ($N=267$, Seed 42)}
\label{tab:benchmark_summary}
\renewcommand{\arraystretch}{0.92}
\resizebox{\textwidth}{!}{%
\begin{tabular}{llcccccc}
\toprule
\textbf{Model / Architecture} & \textbf{Evaluation Paradigm} & \textbf{Params} & \textbf{Precision (\%)} & \textbf{Scam Recall (\%)} & \textbf{Scam-F1 (\%)} & \textbf{Macro-F1 (\%)} & \textbf{Model Size / Format} \\
\midrule
TF-IDF + Linear SVM & Lexical baseline & -- & $88.42 \pm 1.25$ & $79.82 \pm 2.10$ & $83.89 \pm 1.65$ & $88.10 \pm 1.12$ & $12.4\text{ MB}$ (Sparse) \\
Multinomial Naive Bayes & Lexical baseline & -- & $82.15 \pm 1.80$ & $84.20 \pm 2.45$ & $83.16 \pm 2.11$ & $86.54 \pm 1.50$ & $4.2\text{ MB}$ (Count) \\
PhoBERT-base~\cite{nguyen2020phobert} & Baseline B2.1 & 135M & $94.85 \pm 1.30$ & $89.50 \pm 3.85$ & $92.05 \pm 2.65$ & $91.40 \pm 1.80$ & $540.2\text{ MB}$ (FP32) \\
PhoBERT-large~\cite{nguyen2020phobert} & SOTA dense transformer & 370M & $96.10 \pm 0.85$ & $92.40 \pm 2.15$ & $94.20 \pm 1.50$ & $93.85 \pm 1.10$ & $1.48\text{ GB}$ (FP32) \\
FPT ViBERT-base & Baseline B2.2 & 124M & $95.04 \pm 1.11$ & $90.10 \pm 4.12$ & $92.60 \pm 2.84$ & $91.80 \pm 1.95$ & $498.1\text{ MB}$ (FP32) \\
Gemini 3.5 Flash & 5-shot in-context learning & Cloud & $64.66 \pm 3.10$ & $\mathbf{97.33 \pm 1.20}$ & $77.66 \pm 2.25$ & $81.20 \pm 1.85$ & Cloud API ($T=0.0$) \\
\midrule
\textbf{ViSoBERT (Ours, $\mathcal{L}_{\text{WBCE}}$)} & \textbf{Tier 1 Proposed Screener} & \textbf{97.57M} & $89.65 \pm 3.25$ & $95.64 \pm 2.04$ & $92.51 \pm 2.39$ & $94.51 \pm 0.96$ & 98.42 MB (INT8) \\
\bottomrule
\end{tabular}%
}
\par\vspace{6pt}
\parbox{\textwidth}{\footnotesize \textit{Note:} Multi-seed metrics represent mean $\pm$ sample standard deviation across five independent training seeds ($K=5$). Across the five stratified validation partitions ($73, 73, 73, 75, 73$ scam instances respectively), cumulative true positives were $351$ out of $367$ pooled scam instances ($351/367 = 95.64\%$ pooled scam recall, while Gemini 3.5 Flash achieves $365/375 = 97.33\%$). Deterministic evaluation of the released ONNX INT8 model on the locked frozen holdout partition ($N=267$: 192 ham, 75 scam) using the seed 42 checkpoint at operating decision threshold $\theta_{\text{scam}}=0.90$ yields exactly 72/75 = 96.00\% scam recall (exact two-sided 95\% Clopper-Pearson CI: [88.75\%, 99.17\%], $8/192 = 4.17\%$ FPR, $95.88\%$ overall accuracy) with zero sample exclusions. Macro-F1 denotes the unweighted Macro-F1 = $94.51\% \pm 0.96\%$.}
\end{table*}

Under cost-sensitive fine-tuning ($\alpha=5.0$), ViSoBERT achieves $95.64\%$ mean validation scam recall (SD $2.04$ percentage points; cumulative pooled recall $351/367 = 95.64\%$), outperforming PhoBERT-base ($89.50\%$), FPT ViBERT ($90.10\%$), and PhoBERT-large ($92.40\%$). On the locked frozen holdout ($N=267$), INT8 ViSoBERT achieves $96.00\%$ scam recall ($72/75$, exact $95\%$ Clopper-Pearson CI: $[88.75\%, 99.17\%]$) at $4.17\%$ FPR ($8/192$).

In contrast, while solitary Gemini 3.5 Flash achieved $97.33\%$ recall, it exhibited extreme hyper-vigilance, recording $41$ false positives on $192$ benign messages (empirical False Positive Rate of $21.35\%$, $41/192$). Among all 116 scam alerts fired by Gemini, 41 were false alarms, representing an empirical False Discovery Rate of $35.34\%$ ($41/116$). Specifically, the False Discovery Rate was $35.34\%$.

\begin{table}[h]
\centering
\caption{Paired McNemar Test Analysis across Model Configurations ($\alpha=0.05$)}
\label{tab:mcnemar}
\renewcommand{\arraystretch}{0.92}
\resizebox{\columnwidth}{!}{%
\begin{tabular}{lcccccc}
\toprule
\textbf{Paired Model Comparison} & $b$ & $c$ & $\chi^2$ & $p$-value & Odds Ratio & Signif. \\
\midrule
ViSoBERT vs. Gemini 3.5 Flash & 33 & 3 & 23.3611 & $< 0.0001$ & 11.00 & True \\
ViSoBERT vs. PhoBERT-base     & 10 & 5 & 1.07 & 0.4545 & 2.00 & False \\
ViSoBERT vs. PhoBERT-large    & 7  & 7 & 0.00 & 1.0000 & 1.00 & False \\
ViSoBERT vs. FPT ViBERT-base  & 11 & 7 & 1.57 & 0.4807 & 1.57 & False \\
\bottomrule
\end{tabular}%
}
\end{table}

Table~\ref{tab:mcnemar} details paired McNemar tests with Edwards' continuity correction. Comparing ViSoBERT and Gemini 3.5 Flash, ViSoBERT was correct while Gemini erred in $b = 33$ cases, whereas Gemini was correct while ViSoBERT erred in only $c = 3$ cases ($b=33, c=3$, $b - c = 30$, $\chi^2 = 23.3611, p < 0.0001, \text{OR} = 11.00$). The paired difference Newcombe $95\%$ CI: $[-9.75\%, -0.45\%]$ confirms significant superiority over solitary cloud LLMs. For transformer comparisons (PhoBERT-base $p=0.4545$, PhoBERT-large $p=1.0000$, FPT ViBERT $p=0.4807$), no statistically significant difference in paired overall correctness was detected on this holdout, demonstrating that INT8 quantization ($98.42\text{ MB}$, $70.7\text{ ms}$) incurs no detectable accuracy loss relative to dense FP32 baselines.

\subsection{RQ2: Multi-Agent Forensic Council Ablation}
\label{subsec:res_rq2}

Table~\ref{tab:ablation_council} assesses council configurations on the Hard Forensic Challenge Suite ($N=117$: $98$ scam lures, $19$ authentic transactional notices).

\begin{table*}[t]
\centering
\caption{Empirical Performance, Per-Class Metrics, and Cost Matrix on Hard Forensic Challenge Suite ($N=117$, 98 Scam, 19 Benign)}
\label{tab:ablation_council}
\renewcommand{\arraystretch}{0.92}
\resizebox{\textwidth}{!}{%
\begin{tabular}{lcccccccc}
\toprule
\textbf{Experimental Configuration / Arm} & \textbf{Seats} & \textbf{Accuracy (\%)} & \textbf{Scam Recall (\%)} & \textbf{Benign Recall (\%)} & \textbf{Ham Prec. (\%)} & \textbf{Macro-F1 (\%)} & \textbf{Est. Cost / 1k SMS} & \textbf{API Calls / Sample} \\
\midrule
Arm 1: Full 5-Agent Council Specification & 5 & $94.87\%$ & $93.88\%$ & $\mathbf{100.00\%}$ & $76.00\%$ & $90.28\%$ & $\$0.7500$ & $5.00$ ($585$ total) \\
Arm 2: Single-LLM Baseline ($\text{P-single}$) & 1 & $\mathbf{96.58\%}$ & $\mathbf{97.96\%}$ & $89.47\%$ & $89.47\%$ & $\mathbf{94.02\%}$ & $\mathbf{\$0.0450}$ & $\mathbf{1.00}$ ($117$ total) \\
Arm 3: 3-Agent (No Defender, 2 seats) & 2 & $91.45\%$ & $89.80\%$ & $\mathbf{100.00\%}$ & $65.52\%$ & $84.84\%$ & $\$0.4500$ & $3.00$ ($351$ total) \\
Arm 4: Proposed 3-Agent Council (With Defender Veto) & 3 & $\mathbf{90.60\%}$ & $88.78\%$ & $\mathbf{100.00\%}$ & $\mathbf{76.00\%}$ & $86.58\%$ & $\$0.4500$ & $3.00$ ($351$ total) \\
\bottomrule
\end{tabular}%
}
\par\vspace{6pt}
\parbox{\textwidth}{\footnotesize \textit{Note:} Reported metrics reflect evaluations on the 117-sample hard challenge suite. In full-denominator fail-safe evaluation, UNKNOWN abstentions ($0.40 \le P < 0.60$) are mapped to non-scam Ham ($100\%$ coverage), yielding $19/(19 + 2 + 4) = 19/25 = 76.00\%$ Ham precision for Arm 1 and $19/(19 + 1 + 5) = 19/25 = 76.00\%$ for Arm 4. Under selective evaluation (excluding UNKNOWN verdicts), Ham precision is $19/(19+2) = 90.48\%$ for Arm 1 and $18/(18+1) = 94.74\%$ for Arm 4 ($96.58\%$ coverage). Arm 4 (Proposed 3-Agent Council with Defender Veto) raw confusion distribution across 98 scams is 88 SCAM-RISK, 1 LOW-RISK, 9 UNKNOWN; across 19 benign messages, 18 LOW-RISK, 1 UNKNOWN, 0 SCAM-RISK (Sum = 106/117 = 90.60\% overall accuracy). The Clopper-Pearson 95\% CI for 0/19 benign false alarms is [0.00\%, 17.65\%]. Costs are estimated based on $\$0.15/1\text{k}$ unit calls.}
\end{table*}

Ablation reveals critical operational trade-offs:
\begin{enumerate}
    \item \textbf{Single-LLM vs. Multi-Agent Council:} The Single-LLM baseline ($\text{P-single}$, Arm 2) achieved higher raw accuracy ($96.58\%$) and recall ($97.96\%$) at lower estimated cost ($\$0.0450/1\text{k}$). However, P-single recorded $2$ false alarms on authentic notices ($10.53\%$ FPR, $2/19$). In contrast, the proposed 3-agent council with Defender Veto (Arm 4) achieved $0.00\%$ false alarms ($0/19$, exact 95\% Clopper-Pearson CI: $[0.00\%, 17.65\%]$) by delegating 10 ambiguous edge cases to UNKNOWN abstention.
    \item \textbf{3-Agent vs. 5-Agent Council:} Exact two-sided binomial tests detected no statistically significant difference in paired classification correctness between the 3-agent and 5-agent configurations ($p = 0.0625$, Edwards' corrected $\chi^2 = 3.2000$). The 3-agent panel reduced API queries by $40.0\%$ ($585 \to 351$), saving $\$0.30/1\text{k}$ in estimated API fees while scam recall decreased by $5.10$ percentage points ($93.88\% \to 88.78\%$).
    \item \textbf{Public Defender Impact:} Omitting Seat C (Arm 3) caused Ham precision to collapse to $65.52\%$, whereas integrating Seat C with $\delta_{\text{defender}} = 0.25$ (Arm 4) restored Ham precision to $76.00\%$ ($94.74\%$ selective precision). We note that Arm 3 reduces the council size from three agents to two seats rather than substituting a generic control agent, introducing an empirical confounder.
\end{enumerate}

Complete $2 \times 3$ confusion matrices across all four experimental arms are cataloged in Supplementary Table~S3.

\subsection{RQ3: Two-Tier System Efficiency and Cost Trade-offs}
\label{subsec:res_rq3}

Table~\ref{tab:e2e_benchmark} evaluates the end-to-end cascaded architecture on the combined exploratory development benchmark ($N=384$, $\tau=0.90$).

\begin{table*}[t]
\centering
\caption{Comparative End-to-End System Benchmark Across Architectural Arms ($N=384$ Combined Development Samples, $\tau=0.90$)}
\label{tab:e2e_benchmark}
\renewcommand{\arraystretch}{0.92}
\resizebox{\textwidth}{!}{%
\begin{tabular}{lccccccc}
\toprule
\textbf{Architectural Arm} & \textbf{Accuracy (\%)} & \textbf{Precision (\%)} & \textbf{Scam Recall (\%)} & \textbf{FPR (\%)} & \textbf{Mean Latency (ms)} & \textbf{Est. Cost / 1k SMS} & \textbf{Escalation Rate (\%)} \\
\midrule
Arm 1: Solo Tier 1 Edge (ViSoBERT) & $96.35\%$ & $96.49\%$ & $95.38\%$ & $2.84\%$ & $\mathbf{70.7\text{ ms}}$ & $\mathbf{\$0.00}$ & $0.0\%$ \\
Arm 2: Solo Tier 2 Fast (Gemini Flash-Lite) & $97.66\%$ & $95.56\%$ & $99.42\%$ & $3.79\%$ & $1,969.9\text{ ms}$ & $\$0.2817$ & $100.0\%$ \\
Arm 3: Solo Tier 2 SOTA (GPT-6.1-sol) & $99.74\%$ & $100.00\%$ & $99.42\%$ & $\mathbf{0.00\%}$ & $5,765.0\text{ ms}$ & $\$1.4667$ & $100.0\%$ \\
Arm 4: Naive Cascaded (Tier 1 + PoLL 3/5) & $97.14\%$ & $96.55\%$ & $97.11\%$ & $2.84\%$ & $717.8\text{ ms}$ & $\mathbf{\$0.0541}$ & $3.4\%$ \\
\textbf{Arm 5: Proposed ScamShield (Two-Tier)} & $\mathbf{97.40\%}$ & $\mathbf{97.11\%}$ & $\mathbf{97.11\%}$ & $\mathbf{2.37\%}$ & $\mathbf{267.5\text{ ms}}$ & $\mathbf{\$0.0153}$ & $\mathbf{3.4\%}$ \\
\bottomrule
\end{tabular}%
}
\par\vspace{6pt}
\parbox{\textwidth}{\footnotesize \textit{Note:} Reported metrics reflect evaluations on the $N=384$ exploratory development suite. On the frozen holdout ($N=267$), ViSoBERT at default threshold $\theta = 0.50$ records $13$ FP out of $192$ benign ($6.77\%$ FPR) and $74$ TP out of $75$ scam ($98.67\%$ recall); tuning to $\theta_{\text{scam}} = 0.90$ monotonically reduces false positives to $8$ FP out of $192$ ($4.17\%$ FPR) and $72$ TP out of $75$ ($96.00\%$ recall). Across the full $N=384$ exploratory suite, Arm 1 records $6$ FP out of $211$ benign ($2.84\%$ FPR). Arm 4 naive cascaded cost is $\$0.0541$ (scaled at $3.4\%$ escalation). Proposed Arm 5 marginal cloud cost is calculated from the $3.4\%$ escalation rate to the 3-agent council ($\$0.4500 \times 0.034 = \$0.0153$ per 1,000 messages); system latency is $267.5\text{ ms}$ mean system latency (median P50: $70.7\text{ ms}$, P95: $80.5\text{ ms}$, P99: $\approx 5,850\text{ ms}$).}
\end{table*}

At routing threshold $\tau = 0.90$, the edge screener resolves $96.6\%$ of traffic locally at $70.7\text{ ms}$ median latency. Only $3.4\%$ of payloads escalate to Tier 2, achieving $97.40\%$ overall accuracy, $97.11\%$ scam recall, and suppressing FPR to $2.37\%$. Estimated marginal cloud cost is $\$0.0153$ per 1,000 messages, representing a $98.9\%$ cost reduction relative to querying solo GPT-6.1-sol ($\$1.4667/1\text{k}$). Mean system latency across all routed traffic is $267.5\text{ ms}$.

\textbf{Threshold Sensitivity Sweeps:} As detailed in Supplementary Table~S4, varying $\tau \in [0.70, 0.95]$ demonstrates continuous trade-offs: at $\tau = 0.70$, escalation drops to $0.8\%$ ($99.5\text{ ms}$ mean latency, $\$0.0036/1\text{k}$ cost); at $\tau = 0.90$, escalation is $3.4\%$ ($267.5\text{ ms}$, $\$0.0153/1\text{k}$); at $\tau = 0.95$, escalation increases to $7.6\%$ ($540.2\text{ ms}$, $\$0.0342/1\text{k}$). Operating at $\tau = 0.90$ provides an optimal balance between edge autonomy and threat coverage.

\subsection{Threats to Validity and Limitations}
\label{subsec:threats}

Several experimental limitations must be acknowledged:
\begin{enumerate}
    \item \textbf{Data Overlap in Challenge Suite:} The Hard Challenge Suite ($N=117$) shares 10 exact strings with the frozen test partition and contains 99 payloads derived from training templates, making it an exploratory development evaluation rather than a fully disjoint test set.
    \item \textbf{Small Benign Sample Bounds:} The benign challenge subset contains $N=19$ messages. While zero false alarms were observed ($0/19$), the exact two-sided 95\% Clopper-Pearson confidence interval spans $[0.00\%, 17.65\%]$, precluding unhedged claims in unrestricted deployment. Client-side PII sanitization concordance across 117 challenge samples achieved $[0.00\%, 3.10\%]$ bound.
    \item \textbf{Proxy Gateway Execution:} Cloud LLM evaluations utilized proxy gateway endpoints (\texttt{api.vnaipro.app} / \texttt{vibi.top}) with snapshot alias \texttt{gpt-6.1-sol-2026-03}; unverified backend load balancing remains an unmitigated threat.
    \item \textbf{Ablation Scope:} A combined edge + P-single cascade under identical routing was not evaluated on all threshold sweeps; council benefits are documented strictly relative to observed trade-offs.
\end{enumerate}